%%%%%%%%%%%%%%%%%%%%%%%%%%%%%%%%%%%%%%%%%%%%%%%%%%%%%%%%%%%%%%%%%%%%%%%
% BAB 1 - PENDAHULUAN
%%%%%%%%%%%%%%%%%%%%%%%%%%%%%%%%%%%%%%%%%%%%%%%%%%%%%%%%%%%%%%%%%%%%%%%

\mychapter{1}{BAB 1 PENDAHULUAN}

\section{Latar Belakang}

% Tulis latar belakang penelitian di sini.
% Contoh sitasi: \parencite{contoh-artikel}
% Contoh sitasi dalam kalimat: \textcite{contoh-buku}

Lorem ipsum dolor sit amet, consectetur adipiscing elit.

\section{Rumusan Masalah}

% Tulis rumusan masalah di sini.

Berdasarkan latar belakang di atas, rumusan masalah penelitian ini adalah:
\begin{enumerate}
  \item Pertanyaan penelitian pertama?
  \item Pertanyaan penelitian kedua?
\end{enumerate}

\section{Tujuan Penelitian}

% Tulis tujuan penelitian di sini.

Tujuan dari penelitian ini adalah:
\begin{enumerate}
  \item Tujuan pertama.
  \item Tujuan kedua.
\end{enumerate}

\section{Manfaat Penelitian}

% Tulis manfaat penelitian di sini.

\section{Batasan Masalah}

% Tulis batasan masalah di sini.

\section{Sistematika Pembahasan}

\noindent
\textbf{BAB I : PENDAHULUAN}

Bab ini berisi latar belakang, rumusan masalah, tujuan, manfaat,
batasan masalah, dan sistematika pembahasan.

\noindent
\textbf{BAB II : LANDASAN TEORI}

Bab ini berisi tinjauan pustaka dan landasan teori yang mendukung
penelitian.

\noindent
\textbf{BAB III : METODOLOGI}

Bab ini menjelaskan metodologi yang digunakan dalam penelitian.

\noindent
\textbf{BAB IV : HASIL DAN PEMBAHASAN}

Bab ini memaparkan hasil penelitian beserta pembahasannya.

\noindent
\textbf{BAB V : KESIMPULAN DAN SARAN}

Bab ini berisi kesimpulan dan saran berdasarkan hasil penelitian.
